  \documentclass[xcolor=table]{beamer}
\usepackage{beamerthemesplit}
\usepackage{wrapfig}
\usetheme{SPbGU}
\usepackage{pdfpages}
\usepackage{amsmath}
\usepackage{cmap} 
\usepackage[T2A]{fontenc} 
\usepackage[utf8]{inputenc}
\usepackage[english]{babel}
\usepackage{indentfirst}
\usepackage{amsmath}
\usepackage{tikz}
\usepackage{multirow}
\usepackage[noend]{algpseudocode}
\usepackage{algorithm}
\usepackage{algorithmicx}
\usepackage{fancyvrb}
\usetikzlibrary{shapes,arrows}
%usepackage{fancyvrb}
%\usepackage{minted}
%\usepackage{verbments}


\beamertemplatenavigationsymbolsempty

\title[Результаты группы за 2017 год]{Результаты группы за 2017 год}
\institute[СПбГУ]{
JetBrains Research, Programming Languages and Tools Lab  \\
Сенкт-Петербургский Государственный Университет
}

\author[Семён Григорьев]{Семён Григорьев}

\date{16.12.2017}

\definecolor{orange}{RGB}{179,36,31}

\begin{document}
{
\begin{frame}[fragile]
  \begin{tabular}{p{2.0cm} p{7.5cm} p{1cm}}
   \begin{center}
      \includegraphics[height=1.5cm]{pictures/jetbrainsResearch.pdf}
    \end{center}
    &
    \begin{center}
      \includegraphics[height=1.5cm]{pictures/YC_logo.pdf}
    \end{center}
    &
    \begin{center}
      \includegraphics[height=1.5cm]{pictures/SPbGU_Logo.png}
    \end{center} 
  \end{tabular}
  \titlepage
\end{frame}
}


\begin{frame}[fragile]
  \transwipe[direction=90]
  \frametitle{Группы}
Направления работ
\begin{itemize}
      \item Применение методов синтаксического анализа к обработке биологических данных
      \item Задачи о пересечении для различных классов языков: теория и практика
      \item Интерливинг (shuffle) языков: теория и практика 
\end{itemize}

Состав
\begin{itemize}
      \item Магистры (грант): 2
      \item Бакалавры 3 курс (стипендия): 2
      \item Студенты СПбГУ, ЛЭТИ, АУ: 7
\end{itemize}
\end{frame}


\begin{frame}[fragile]
  \transwipe[direction=90]
  \frametitle{Конференции}
\begin{itemize}
      \item \textbf{TMPA-2017}  Артём Горохов, доклад
            \begin{itemize}
      \item Extended Context-Free Grammars Parsing with Generalized LL
      \end{itemize}

      \item \textbf{PLC-2017}   Дмитрий Ковалёв, Рустам Азимов, доклад
            \begin{itemize}
      \item Динамически формируемый код: синтаксический анализ контекстно-свободной аппроксимации
      \item Синтаксический анализ графов и задача генерации строк с ограничениями
      \end{itemize}

      \item \textbf{SECR-2017}  Семён Григорьев, доклад
            \begin{itemize}
      \item Context-Free Path Querying with Structural Representation of Result
      \end{itemize}

      \item \textbf{BIATA-2017} Семён Григорьев, постер
      \begin{itemize}
      \item Graph parsing application for bioinformatics problems
      \end{itemize}
      \item \textbf{Parsing@SLE-2017} Екатерина Вербицкая, доклад
            \begin{itemize}
      \item Parsing techniques for graph analysis
      \end{itemize}

\end{itemize}
\end{frame}

\begin{frame}[fragile]
  \transwipe[direction=90]
  \frametitle{Публикации}
\begin{itemize}
      \item Материалы конференций
        \begin{itemize}
          \item \textbf{TMPA-2017} --- springer
          \item \textbf{PLC-2017} --- ВАК
          \item \textbf{SECR-2017} --- ACM
        \end{itemize}
      \item Azimov R., Grigorev S. Graph Parsing by Matrix Multiplication //arXiv preprint arXiv:1707.01007. – 2017.
\end{itemize}
\end{frame}

\begin{frame}[fragile]
  \transwipe[direction=90]
  \frametitle{Сотрудничество}
\begin{itemize}
      \item Опубликована информация о стажировках
      \item Заявка на грант РНФ совместно с Александром Охотиным
      \item Семинар по формальным языкам
      \item Prof. Martin A. Musicante, Federal Univ. of Rio Grande do Norte
\end{itemize}
\end{frame}

\end{document}
